\section*{अध्याय \ref{ch:g1:animals-around-us} --- हमारे आस-पास के जंतु}

\begin{solution}{exo:g1:animals-around-us:1}
\omterm{def:g1:animals-around-us:parts}{सिर}, धड़, टाँगें (\omterm{def:g1:animals-around-us:parts}{उपांग}), पूँछ। \omterm{def:g1:animals-around-us:parts}{सिर} आँखें और मुँह लिए रहता है।
\end{solution}

\begin{solution}{exo:g1:animals-around-us:2}
(क) $2$; (ख) $4$; (ग) $6$; (घ) $8$; (ङ) एक भी नहीं --- केंचुए की कोई
टाँग नहीं होती।
\end{solution}

\begin{solution}{exo:g1:animals-around-us:3}
भेड़: ऊन, जो मोटा \omterm{def:g1:animals-around-us:covering}{रोम} है। बत्तख़: \omterm{def:g1:animals-around-us:covering}{पंख}। ट्राउट: \omterm{def:g1:animals-around-us:covering}{शल्क}। घोंघा: एक \omterm{def:g1:animals-around-us:covering}{खोल}
(मुलायम, नम शरीर के ऊपर)।
\end{solution}

\begin{solution}{exo:g1:animals-around-us:4}
कलचिड़ी, मधुमक्खी और तितली। इन सबके डैने होते हैं।
\end{solution}

\begin{solution}{exo:g1:animals-around-us:5}
घोंघा अपने एक मुलायम पैर पर धीरे-धीरे सरकता है और पीछे लुआब की लकीर
छोड़ जाता है। केंचुआ बल खाता है: वह लंबा और पतला होता है, फिर अपने आप
को छोटा कर लेता है, और थोड़ा-थोड़ा आगे बढ़ता जाता है।
\end{solution}

\begin{solution}{exo:g1:animals-around-us:6}
वह अपनी टाँगों पर चलती है, अपने झिल्लीदार पैरों से तैरती है (वही
टाँगें, जिनकी उँगलियाँ \omterm{def:g1:animals-around-us:covering}{त्वचा} से जुड़ी होती हैं), और अपने डैनों से
उड़ती है।
\end{solution}

\begin{solution}{exo:g1:animals-around-us:7}
चींटी दूसरे सजीवों को खाकर पोषण लेती है (भोजन के टुकड़े, बीज, नन्हे
\omterm{def:g1:animals-around-us:animal}{जंतु}), और वह अपनी छह टाँगों पर ख़ुद इधर-उधर जाती है। पोषण और चाल:
परिभाषा के दोनों हिस्से बैठते हैं।
\end{solution}

\begin{solution}{exo:g1:animals-around-us:8}
उत्तर खुला है। अच्छे वर्णन में दिखे हुए भागों के नाम, टाँगों की सही
गिनती, आवरण, चलने का ढंग और जगह होती है --- जैसे: ``दीवार पर एक
घोंघा: दो स्पर्शकों वाला \omterm{def:g1:animals-around-us:parts}{सिर}, मुलायम धड़, कोई टाँग नहीं, भूरा \omterm{def:g1:animals-around-us:covering}{खोल},
वह सरकता है; बारिश के बाद मिला।''
\end{solution}

\begin{solution}{exo:g1:animals-around-us:9}
चलना चलने का सिर्फ़ एक ढंग है। मछलियों के मीन-पंख होते हैं और वे
तैरती हैं --- यही उनका ढंग है। \omterm{def:g1:animals-around-us:animal}{जंतु} को पोषण लेना और चलना चाहिए, चलना
पैरों से ही हो, यह ज़रूरी नहीं: मछली दोनों करती है, इसलिए वह \omterm{def:g1:animals-around-us:animal}{जंतु}
है।
\end{solution}

\begin{solution}{exo:g1:animals-around-us:10}
हाँ, वह पक्षी है: \omterm{def:g1:animals-around-us:covering}{पंख} और डैने ही भेद खोल देते हैं, भले ही वह उड़ता
न हो। शुतुरमुर्ग़ उड़ नहीं सकता पर किसी भी आदमी से तेज़ दौड़ता है ---
पक्षी होना शरीर की बात है, उड़ान की नहीं।
\end{solution}
